\chapter[LA PHILOSOPHIE DE CONCEPTION APRS]{La philosophie de conception
APRS}\label{la-philosophie-de-conception-aprs}

\begin{aprssideblock}{Net Cycle Time}
\phantomsection\addcontentsline{toc}{section}{Net Cycle Time}\label{net-cycle-time}

APRS est avant tout un outil de communication tactique temps réel,
destiné à faciliter le flux d'information lors d'événements spéciaux,
urgences, Skywarn, centre d'opérations d'urgence, ou simplement usage
terrain sous stress. Mais comme dans le monde réel, 99 \% du temps il
fonctionne en routine, en attente de l'événement sérieux improbable.

Toute évolution d'APRS doit préserver sa capacité à fonctionner
localement sous stress. Détails de cette philosophie :

\begin{enumerate}
\def\labelenumi{\arabic{enumi}.}
\item
  APRS utilise le concept de « net cycle time » --- durée dans laquelle
  un utilisateur doit avoir entendu (au moins une fois) toutes les
  stations APRS à portée, pour obtenir une image plus ou moins complète
  de l'activité. Le net cycle time varie avec les conditions locales et
  le nombre de digipeaters traversés.
\item
  Objectif : 10 minutes de net cycle time en usage local. En 10 minutes
  après l'arrivée sur les lieux, on doit avoir capturé toute la
  situation tactique.
\item
  Toutes les stations, même fixes, doivent beaconer leur position au
  rythme du net cycle time. En situation de stress, les stations vont et
  viennent en permanence. Les rapports de position montrent où elles
  sont sans qu'on ait à demander, et signalent qu'elles sont toujours
  actives.
\item
  On ne peut pas supposer que tous les utilisateurs APRS répondant à un
  événement comprennent les ramifications d'APRS et les statistiques du
  canal --- les réglages utilisateurs ne suffisent pas pour éviter de
  tuer un réseau sous stress. Pour anticiper une saturation, APRS ajuste
  automatiquement son net cycle time selon le nombre de digipeaters dans
  le chemin UNPROTO :

  \begin{itemize}
  \tightlist
  \item
    Opération directe (aucun digipeater) : 10 minutes (probablement un
    événement).
  \item
    1 hop : 10 minutes (probablement un événement).
  \item
    2 hops : 20 minutes.
  \item
    3 hops ou plus : 30 minutes.
  \end{itemize}
\item
  Comme quasiment toutes les stations fixes règlent leur chemin sur 3
  digis ou plus, le net cycle time par défaut en usage quotidien routine
  est de 30 minutes. Standard universel : allume la radio et APRS, ne
  fais rien d'autre, en 30 minutes tu as l'image complète de toutes les
  stations APRS à portée.
\item
  Comme la connaissance de la localisation des digipeaters est
  fondamentale, les digipeaters devraient utiliser plusieurs commandes
  de beacon pour transmettre à différents rythmes et sur différents
  chemins --- par ex. 10 min localement, 30 min via 3 digis (et d'autres
  selon besoin).
\item
  Si le net cycle time est trop long, les utilisateurs seront tentés
  d'envoyer des requêtes de stations, ce qui augmentera inutilement le
  trafic. Extrêmes recommandés : 10 et 30 minutes --- ça donne aux
  concepteurs de réseau les hypothèses de charge canal nécessaires à une
  bonne ingénierie.
\end{enumerate}
\end{aprssideblock}

\begin{aprssideblock}{Packet Timing}
\phantomsection\addcontentsline{toc}{section}{Packet Timing}\label{packet-timing}

Les paquets APRS sont sans erreur (grâce à la FCS AX.25) mais sans
livraison garantie. APRS transmet donc l'information de façon
redondante. Pour assurer une livraison rapide des données nouvelles ou
changeantes, et préserver la capacité canal en réduisant les
interférences avec les données anciennes, APRS doit transmettre les
données récentes plus souvent que les anciennes.

Plusieurs algorithmes en usage :

\begin{itemize}
\tightlist
\item
  \textbf{Decay Algorithm} --- transmettre un nouveau paquet une fois,
  puis n secondes plus tard. Doubler n à chaque nouvelle transmission.
  Quand n atteint le net cycle time, rester à ce rythme. D'autres
  facteurs que « doubler » sont possibles selon le contexte.
\item
  \textbf{Fixed Rate} --- transmettre une fois, puis n secondes plus
  tard. Répéter x fois puis s'arrêter.
\item
  \textbf{Message-on-Heard} --- transmettre selon l'un des algorithmes
  ci-dessus. Si le paquet est toujours valide et n'a pas été acquitté,
  et si le net cycle time est atteint, le destinataire est probablement
  absent. Cependant si un paquet du destinataire est ensuite entendu,
  réessayer une fois de transmettre.
\item
  \textbf{Time-Out} --- période au-delà de laquelle on peut
  raisonnablement considérer qu'une station n'existe plus ou est off-air
  si aucun paquet n'a été entendu. Défaut nominal : 2 heures. Ce
  time-out n'est pas utilisé dans les algorithmes de transmission mais
  sert à décider quand cesser d'afficher une station comme « active ».
  Sur HF, les signaux vont et viennent --- décisions plus flexibles.
\end{itemize}
\end{aprssideblock}

\begin{aprssideblock}{Digipeating\\générique}
\phantomsection\addcontentsline{toc}{section}{Digipeating générique}\label{digipeating-guxe9nuxe9rique}

La puissance d'APRS sur le terrain vient de l'usage du digipeating
générique : les paquets se propagent sans connaissance a priori du
réseau. Six techniques ont évolué depuis 1992 :

\begin{enumerate}
\def\labelenumi{\arabic{enumi}.}
\tightlist
\item
  \textbf{RELAY} --- tout TNC APRS VHF est supposé avoir l'alias
  \texttt{RELAY}, utilisable comme digipeater par n'importe qui.
\item
  \textbf{ECHO} --- les stations HF utilisent l'alias \texttt{ECHO}
  comme alternative à \texttt{RELAY}. Attention à la propagation HF qui
  peut causer des interférences sur une vaste zone --- usage
  parcimonieux.
\item
  \textbf{WIDE} --- tout digipeater en hauteur est supposé avoir l'alias
  \texttt{WIDE} pour couvrir plus loin.
\item
  \textbf{TRACE} --- tout digipeater en hauteur pratiquant la
  substitution d'indicatif est supposé avoir l'alias \texttt{TRACE}. Ces
  digipeaters s'auto-identifient dans les paquets qu'ils digipeatent en
  substituant leur indicatif à \texttt{RELAY}, \texttt{WIDE} ou
  \texttt{TRACE}.
\item
  \textbf{WIDEn-N} --- un digipeater supportant WIDEn-N digipeatera tout
  paquet WIDEn-N « nouveau » et décrémentera le SSID jusqu'à ce qu'il
  atteigne \texttt{-0}. Le digipeater garde une copie ou un checksum et
  ne re-digipeatera pas le même paquet dans un intervalle typique de 28
  s. Réduit considérablement les doublons dans les zones où plusieurs
  digipeaters s'entendent.
\item
  \textbf{GATE} --- indicatif générique pour les gateways HF-VHF. Tout
  paquet entendu en HF via \texttt{GATE} est digipeaté localement en
  VHF. Les réseaux locaux gardent un œil sur l'échelle
  nationale/internationale.
\end{enumerate}
\end{aprssideblock}

\begin{aprssideblock}{Vues\\cartographiques}
\phantomsection\addcontentsline{toc}{section}{Communiquer des vues cartographiques sans ambiguïté}\label{communiquer-des-vues-cartographiques-sans-ambiguuxeftuxe9}

APRS est un système géographique tactique. Pour maximiser son efficacité
opérationnelle et minimiser la confusion entre opérateurs, il faut une
manière non ambiguë de communiquer la « localisation » et la « taille »
(aire de couverture) d'une vue.

La convention APRS : référence à un \emph{centre} et à un \emph{range}
qui spécifient le centre géographique et le rayon approximatif d'un
cercle inscrit dans la vue, indépendant de l'aspect ratio. Ce rayon (en
milles nautiques, milles terrestres ou km) est le « range scale ». Base
commune simple pour décrire n'importe quelle vue à d'autres, quel que
soit le médium ou le programme.
\end{aprssideblock}
